\section{讨论：构型优势、局限与扩展路径}
\label{sec:discussion}

\subsection{构型核心优势}

\begin{enumerate}[label=(\arabic*)]
    \item \textbf{主通道结构清晰}：正交摆轴使推力力臂主项形成对角结构，
          可用低复杂度直接反馈或近似分配；完整效能矩阵仍含反扭矩交叉项，且随推力下降而退化；
    \item \textbf{形态独立的低速控制来源}：在推进系统保持非零推力且执行器未饱和时，
          摆动力矩不以自由来流动压为前提，可在有翼、无翼和低速形态下提供姿态控制。
          当前Web core已回归验证局部悬停平衡与小扰动恢复；垂起过渡图仍来自历史脚本，
          尚未按当前参数、轴置换和固件限幅重跑，也未验证有翼/无翼构型切换和完整过渡走廊，
          因此不能据此宣称全包线转换能力；
    \item \textbf{反扭主动对消与偏航控制一体化}：前后反向旋转螺旋桨使名义工况
          的轴向反扭矩近似抵消，并允许通过差速主动生成模型$x'$轴力矩。
          经当前固件轴置换后，该通道承担偏航主控；上摆通过推力力臂承担滚转主控。
          这一对应关系不能再按旧巡航观察方向写成“差速滚转”；
    \item \textbf{四输入主通道设计}：两台电机、两个单轴摆座共同提供总推力与三轴主力矩输入。
          这减少了分配维数，但不等同于已经证明全包线可控、故障容错或执行器数量全局最优；
    \item \textbf{尾翼权衡空间}：推力矢量可承担部分俯仰和偏航控制任务，因而允许重新评估尾翼尺寸。
          是否减尾翼必须同时验证静稳定、阻尼、滑翔与失效状态，当前模型不足以证明去尾翼收益。
\end{enumerate}

\subsection{已知局限与设计权衡}

任何构型设计都是在矛盾的需求之间寻找工程上可行的折中。以下列出本构型的主要局限，
它们不是缺陷，而是设计决策的内在权衡——了解这些权衡对于在正确的应用场景中
使用这一构型至关重要。

\begin{enumerate}[label=(\arabic*)]
    \item \textbf{低油门推进姿态效能退化}：差速和两个摆座的名义效能均随
          $\omega_0^2$量级下降。固件任务语义中差速对应偏航、上摆对应滚转、下摆对应俯仰；
          因此最低油门应由三轴控制裕度共同约束。有翼前飞时可由气动舵面补充，悬停形态则需
          维持足够推重比和饱和裕度；
    \item \textbf{$w_{\max}$饱和导致推力偏离}：大差速+高油门组合可能导致
          一侧电机超速。转速硬限幅虽然保护了电机，但破坏了推力不变假设。
          这会在机动最剧烈时产生额外的总推力扰动；
    \item \textbf{俯仰-偏航耦合不可完全消除}：反扭矩的摆角投影耦合
          当前MODEL-DEFAULT下主通道尺度比$k_Q/(k_Tb)\approx8.55\%$，但$k_T$、$k_Q$
          尚无独立台架标定；实机调参只提供量级锚点，实际耦合还取决于角度比和工作点。
          通过前馈补偿可降低耦合，效果需闭环验证；
    \item \textbf{推进故障下完整三轴控制不可保证}：纯推力矢量构型在双发失效后丧失推进控制；
          单发失效也会同时改变总推力、可用摆动通道和滚转反扭矩，不能假定仍保持完整三轴控制。
          发动机空中重启能力和供电系统冗余是安全底线要求；
    \item \textbf{欧拉角表示边界}：3-2-1欧拉角在$\theta=\pm90^\circ$处奇异，
          因而Web旧前飞SAS不能直接用于垂直姿态。当前固件滚转/俯仰外环采用四元数短弧误差，
          偏航采用ratchet heading hold；欧拉角主要保留用于显示与航向参考\cite{markley2014fundamentals}；
    \item \textbf{推力方向余弦损失}：摆角偏转使推力有效前向分量降低$\cos\delta$倍。
          $\delta = 15^\circ$满偏时前向推力损失约$3.4\%$（$\cos 15^\circ \approx 0.966$；早期 25° 口径 10\% 为旧限幅残留），
          该百分比不会随时间继续“累计”，但持续偏转会持续增加能耗并压缩航程或推力裕度，
          因而应在性能计算和任务规划中计入。
\end{enumerate}

\subsection{与现有构型的综合比较}

\begin{table}[H]
\centering
\caption{纵列双发矢量推力构型与现有典型构型的定性比较}
\setlength{\tabcolsep}{3pt}
\small
\begin{tabularx}{\textwidth}{>{\raggedright\arraybackslash}p{2.35cm}*{4}{>{\raggedright\arraybackslash}X}}
\toprule
\textbf{特性} & \textbf{本文模块化构型} & \textbf{常规固定翼} & \textbf{四旋翼} & \textbf{倾转旋翼} \\
\midrule
低速姿态效能来源 & 非零推力下的摆动与反扭差 & 动压驱动的舵面 & 旋翼推力差 & 旋翼推力与倾转机构 \\
巡航效率依据 & 有翼形态可由机翼承担升力；无翼形态以推进升力为主，均尚未标定 & 取决于具体气动设计 & 通常受盘载与姿态限制 & 取决于转换机构与旋翼/机翼折中 \\
分配问题 & 三主通道近似，含交叉项 & 常规舵面混控 & 固定几何混控矩阵 & 随构型变化的分配与调度 \\
机械构成 & 2电机+2单轴摆座 & 电机/发动机+多舵面 & 4电机（标准四旋翼） & 多旋翼+倾转机构 \\
垂直起降 & 无翼/有翼垂起局部已仿真；完整转换未验证 & 通常不具备 & 具备 & 具备 \\
单一执行器失效 & 完整三轴控制不保证 & 取决于舵面与布局 & 标准四旋翼通常不保持完整控制 & 强依赖具体冗余设计 \\
\bottomrule
\end{tabularx}
\end{table}

\begin{figure}[H]
\centering
\resizebox{0.96\linewidth}{!}{%
\begin{tikzpicture}[>=Stealth,
  card/.style={draw=black!60, fill=black!3, rounded corners=3pt, minimum width=38mm, minimum height=26mm, text width=34mm, align=center, font=\footnotesize\sffamily, inner sep=3pt},
  result/.style={draw=purple!65!black, fill=purple!5, rounded corners=3pt, minimum width=92mm, minimum height=23mm, text width=86mm, align=left, font=\footnotesize\sffamily, inner sep=4pt},
  mechanism/.style={draw=black!45, fill=white, rounded corners=2pt, minimum width=28mm, minimum height=7mm, align=center, font=\scriptsize\sffamily, inner sep=2pt},
  arrow/.style={->, draw=black!65, line width=.9pt, >={Stealth[length=2.5mm]}},
  source/.style={font=\scriptsize\sffamily, fill=white, inner sep=1.5pt, align=center},
]
  \node[card, draw=blue!60!black, fill=blue!4] (fixed) at (0,0) {\textbf{常规固定翼}\\[2pt]
    \textbf{核心机制：}气动升力面与舵面\\
    \textbf{优势：}巡航效率高\\
    \textbf{局限：}低速控制弱、依赖跑道};
  \node[card, draw=green!55!black, fill=green!4] (quad) at (4.8,0) {\textbf{四旋翼}\\[2pt]
    \textbf{核心机制：}多电机差速控制\\
    \textbf{优势：}悬停与低速操纵强\\
    \textbf{局限：}高速巡航效率低};
  \node[card, draw=orange!70!black, fill=orange!5] (tilt) at (9.6,0) {\textbf{倾转旋翼}\\[2pt]
    \textbf{核心机制：}推力方向连续倾转\\
    \textbf{优势：}兼顾VTOL与巡航\\
    \textbf{局限：}机构重、过渡控制复杂};

  \node[mechanism, draw=blue!60!black] (lift) at (0,-2.55) {固定翼升力面};
  \node[mechanism, draw=green!55!black] (reaction) at (4.8,-2.55) {反扭矩差动};
  \node[mechanism, draw=orange!70!black] (vector) at (9.6,-2.55) {低速矢推控制};
  \draw[arrow, draw=blue!60!black] (fixed.south) -- (lift.north);
  \draw[arrow, draw=green!55!black] (quad.south) -- (reaction.north);
  \draw[arrow, draw=orange!75!black] (tilt.south) -- (vector.north);

  \coordinate (merge) at (4.8,-3.7);
  \draw[arrow, draw=blue!60!black] (lift.south) |- (merge);
  \draw[arrow, draw=green!55!black] (reaction.south) -- (merge);
  \draw[arrow, draw=orange!75!black] (vector.south) |- (merge);
  \fill[black!65] (merge) circle (2pt);
  \node[source, above=3pt] at (merge) {几何融合与四输入主通道};

  \node[result] (ours) at (4.8,-6.45) {\textbf{本文构型：纵列双发正交单轴矢量推力}\\[3pt]
    \textbf{机制组合：}固定翼巡航升力 + 正交单轴矢推 + 反扭矩差动\\
    \textbf{概念潜力：}总推力、两摆角与差速形成结构清晰的四输入主通道\\
    \textbf{方向约定：}前 CW，$\bm h_f\parallel +x_b$；尾 CCW，$\bm h_t\parallel -x_b$；$\delta_f$ 绕 $z_b$，$\delta_t$ 绕 $y_b$（见图~\ref{fig:thrust_mapping}）\\
    \textbf{验证边界：}交叉项、低油门退化、失效状态及巡航效率尚待标定与试验。};
  \draw[arrow, draw=purple!65!black, line width=1.15pt] (merge) -- (ours.north);
\end{tikzpicture}%
}
\caption{构型的机制来源与概念定位：固定翼升力面、反扭矩差动，以及由两个正交单轴摆座产生的矢量控制。图示表达设计关系，不代表已验证的性能排序或VTOL能力。}
\label{fig:config_comparison}
\end{figure}


\subsection{模块化形态与参数切换}
\label{sec:configuration_variants}

将机翼视为可选模块后，动力学模型不能简单地只切换气动力开关。
有翼和无翼形态至少具有不同的质量、质心、转动惯量、气动面积和气动导数：
\begin{equation}
    \mathcal{P}_{\mathrm{wing}}=
    \{m,I_x,I_y,I_z,x_{\mathrm{CG}},S,\bar c,b,C_{*}\},
    \qquad
    \mathcal{P}_{\mathrm{wingless}}=
    \{m',I_x',I_y',I_z',x_{\mathrm{CG}}',S',\bar c',b',C_{*}'\}.
    \label{eq:configuration_params}
\end{equation}
其中$C_{*}$代表气动导数集合。通常$m'\neq m$、$\mathbf{I}'\neq\mathbf{I}$，
且机翼拆装可能引起质心位置和推进器相对力臂变化。
因此形态切换应同步更新：
\begin{enumerate}[label=(\arabic*)]
    \item 质量、质心和惯量张量；
    \item 推重比、悬停转速和电机饱和裕度；
    \item 气动面积、升阻特性和气动阻尼；
    \item 控制效能矩阵、配平点和外环/内环增益调度表。
\end{enumerate}

本文当前采用同一组\mbox{MODEL-DEFAULT}集总参数开展概念级对比。
因此，第\ref{sec:vtol}章的VTOL悬停分析以及第\ref{sec:cascade_chapter}章的
无翼或竖直姿态结果，只能解释为\textbf{推进与姿态控制机制的局部可行性验证}，
不能替代机翼拆装后的参数重辨识、配平计算和全包线试验。

\subsection{技术成熟度与验证路径}
\label{sec:verification}

本文已完成局部构型机理、MODEL-DEFAULT仿真与软件回归，2026-08起项目记录进入
\textbf{实机初步验证阶段}。姿态控制链与FPV摇杆曲线已部署，并记录了首飞、手动直通和调参结果；
但论文工程尚未纳入带时间戳的原始日志、重复试验和参数快照，故这些记录只支持部署状态与方向语义，
不支持完整性能或包线结论。详见第\ref{sec:cascade_implementation}节。
建议的技术成熟度推进路径为（标注各 Tier 的当前进展）：

\textbf{Tier 1 — 参数标定}：在推力台架上测量$k_T(\omega)$、$k_Q(\omega)$、
$\tau_m(\omega)$函数关系，标定电机-桨-ESC系统模型。同步进行CFD计算，
获取气动导数数据库（$C_L(\alpha,\beta)$、$C_D(\alpha,\beta)$、
$C_m(\alpha,\beta,\delta_t)$等）。

\textbf{Tier 2 — 硬件在环（HIL）}：将标定后的参数模型和SAS控制律部署到
实时仿真平台（如Pixhawk HIL或dSPACE），连接真实的飞控硬件（IMU、舵机、电调），
验证控制律的实时性和执行器接口的正确性。

\textbf{Tier 3 — 系留飞行}：在有安全绳约束的条件下进行有限自由度的飞行试验。
系留绳索用于断电回收，并验证非零空速/非零推力下的基本姿态控制。
\textbf{（2026-08项目记录）}受限飞行、手动直通、角速度方波激励和油门瞬态调参已开展；
结论仍需原始日志与重复试验复核。当前$\pm15^\circ$尾摆限幅下，水平机身时单电机垂向分量比例
约为$\sin15^\circ\approx0.259$；机头朝天悬停的重量由名义推力轴承担，本身不要求摆座转过$90^\circ$。

\textbf{Tier 4 — 自由飞行}：从常规跑道起飞和小扰动平飞开始扩展包线，再逐步验证低速、巡航、
大机动和边界工况。若研究悬停到巡航的连续过渡，应先在现有$\pm15^\circ$限幅内求解可行走廊；
只有约束求解证明不可行时，才需要提出新的摆座或推进布局基线。

每一步应设置明确的通过/不通过判据，并记录与仿真预测的偏差以回溯更新模型。

\subsection{未来扩展方向}

\begin{enumerate}[label=(\arabic*)]
    \item \textbf{控制分配升级}：从直接映射升级为加权伪逆+零空间优化，
          在保持主通道对角解耦的同时，利用剩余自由度实现耦合项前馈补偿
          和执行器负载均衡。特别是当引入气动舵面后，执行空间由3扩至6+，
          控制分配的必要性凸显；
    \item \textbf{过渡走廊研究}：若未来研究悬停到平飞的连续过渡，应先在现有
          $\pm15^\circ$执行器约束内开展可达性、配平族与动态走廊求解；只有约束求解证明不可行时，
          才需要重构摆座或整机方案。该方向不属于当前C0仿真能力；
          过渡段的气动-推力耦合是关键控制问题——
          倾转旋翼的过渡控制已有大量研究，其增益调度策略和模型预测控制（MPC）方法
          可为本构型提供参考\cite{kang2009control,kim2003nonlinear}。对于纵列双发构型，
          可为后续新构型研究提供方法参考；
    \item \textbf{全气动-推力融合}：在中高速巡航段引入气动舵面，
          形成推力矢量（低速不依赖动压）+舵面（巡航控制）的混合控制架构。
          混合控制的关键在于设计合理的控制分配策略，在舵面权限充足时优先使用舵面
          （减小推力损失和摆座磨损），在舵面不足时无缝过渡到推力矢量；
    \item \textbf{非线性控制}：在大姿态变化、显著轴间耦合和执行器接近饱和时，
          固定工作点线性SAS的适用性下降。可考虑引入模型预测控制（MPC）、
          滑模控制或反馈线性化等非线性控制方法，在发挥矢量推力机动潜力的同时
          保证稳定性；
    \item \textbf{故障容错}：设计执行器失效后的控制重构策略。对于电机单发失效，
          可设计降级控制律（利用剩余的单发电机的推力矢量和反扭矩维持部分控制），
          将飞行器引导至安全迫降。对于舵机失效，可设计锁定摆座+备用电机的应急策略。
\end{enumerate}
